% GENERATED FILE. Edit canonical lessons, then run npm run generate:books.
% canonical-source-hash: fnv1a64:2f7807f414cef7e4
% canonical-lessons: IT-C283-redazione, IT-C283-cronista, IT-C283-inviato, IT-C283-intervista, IT-C283-intervistare

\chapter{In redazione}
\label{ch:it-in-redazione}
\hlchaptermodality{283}

\begin{chapteropening}
\textbf{By the end of this chapter:} \emph{I can talk about who makes the news: the newsroom, the reporter, the correspondent and the interview.}

Complete the last lesson of chapter 283: I can talk about who makes the news: the newsroom, the reporter, the correspondent and the interview.
\end{chapteropening}

\section[la redazione]{la redazione — the newsroom, the editorial staff}

\label{lesson:IT-C283-redazione}



\begin{quote}
Before the new one: say the Italian for the headline, then the Italian for the article.
\end{quote}

\subsection*{You'll want to know: la redazione}
\textbf{la redazione} — \textquotedblleft{}the newsroom, the editorial staff\textquotedblright{}.

\begin{quote}
In redazione lavorano venti persone.
\end{quote}

\begin{quote}
\emph{Twenty people work in the newsroom.}
\end{quote}

Also the writing of a text: \emph{la redazione di un documento}.

\begin{cousinweb}
\emph{redazione} is usually traced to French \emph{rédaction}, built on Latin \emph{redāctus}, the past participle of \emph{redigere} \textquotedblleft{}to bring back, to put in order\textquotedblright{} (Lewis \& Short).
\end{cousinweb}

\begin{grammarlens}[title={What you've built}]
One more word for telling, explaining and describing at B1.
\end{grammarlens}

\subsection*{Your turn}
\begin{itemize}
\raggedright
  \item \emph{Say it:} \emph{la redazione}
  \item \emph{Say it:} \emph{la redazione}, once more
  \item \emph{Say it:} a sentence of your own with \emph{la redazione}
  \item \emph{Recall:} say the Italian for the headline, then the Italian for the article, then say \emph{la redazione} again
\end{itemize}

\begin{tcolorbox}[breakable,colback=teal!4,colframe=teal!35!black,title={Before you move on}]
What does la redazione mean? (\textquotedblleft{}The newsroom, the editorial staff\textquotedblright{}.) Say it once more.
\end{tcolorbox}
\section[il cronista]{il cronista — the news reporter}

\label{lesson:IT-C283-cronista}



\begin{quote}
Before the new one: say the Italian for the article, then the Italian for the newsroom.
\end{quote}

\subsection*{You'll want to know: il cronista}
\textbf{il cronista} — \textquotedblleft{}the news reporter\textquotedblright{}.

\begin{quote}
Il cronista ha parlato con i vicini.
\end{quote}

\begin{quote}
\emph{The reporter spoke to the neighbours.}
\end{quote}

A woman is \emph{la cronista}: same form.

\begin{cousinweb}
\emph{cronista} is built on \emph{cronaca} \textquotedblleft{}chronicle, news\textquotedblright{} with the ending \emph{-ista}; the root is Greek \emph{khrónos} \textquotedblleft{}time\textquotedblright{} (EtymDB).
\end{cousinweb}

\begin{grammarlens}[title={What you've built}]
One more word for telling, explaining and describing at B1.
\end{grammarlens}

\subsection*{Your turn}
\begin{itemize}
\raggedright
  \item \emph{Say it:} \emph{il cronista}
  \item \emph{Say it:} \emph{il cronista}, once more
  \item \emph{Say it:} a sentence of your own with \emph{il cronista}
  \item \emph{Recall:} say the Italian for the article, then the Italian for the newsroom, then say \emph{il cronista} again
\end{itemize}

\begin{tcolorbox}[breakable,colback=teal!4,colframe=teal!35!black,title={Before you move on}]
What does il cronista mean? (\textquotedblleft{}The news reporter\textquotedblright{}.) Say it once more.
\end{tcolorbox}
\section[l'inviato]{l'inviato — the correspondent, the reporter sent out}

\label{lesson:IT-C283-inviato}



\begin{quote}
Before the new one: say the Italian for the newsroom, then the Italian for the news reporter.
\end{quote}

\subsection*{You'll want to know: l'inviato}
\textbf{l'inviato} — \textquotedblleft{}the correspondent, the reporter sent out\textquotedblright{}.

\begin{quote}
L'inviato del giornale è arrivato ieri.
\end{quote}

\begin{quote}
\emph{The paper's correspondent arrived yesterday.}
\end{quote}

The past participle of \emph{inviare}, ‘to send’; \emph{un inviato speciale} is a special correspondent.

\begin{cousinweb}
\emph{inviato} is the past participle of \emph{inviare} \textquotedblleft{}to send\textquotedblright{}, usually traced to Late Latin \emph{inviāre} \textquotedblleft{}to put on the way\textquotedblright{}, built on \emph{via} \textquotedblleft{}way\textquotedblright{} (Lewis \& Short).
\end{cousinweb}

\begin{grammarlens}[title={What you've built}]
One more word for telling, explaining and describing at B1.
\end{grammarlens}

\subsection*{Your turn}
\begin{itemize}
\raggedright
  \item \emph{Say it:} \emph{l'inviato}
  \item \emph{Say it:} \emph{l'inviato}, once more
  \item \emph{Say it:} a sentence of your own with \emph{l'inviato}
  \item \emph{Recall:} say the Italian for the newsroom, then the Italian for the news reporter, then say \emph{l'inviato} again
\end{itemize}

\begin{tcolorbox}[breakable,colback=teal!4,colframe=teal!35!black,title={Before you move on}]
What does l'inviato mean? (\textquotedblleft{}The correspondent, the reporter sent out\textquotedblright{}.) Say it once more.
\end{tcolorbox}
\section[l'intervista]{l'intervista — the interview}

\label{lesson:IT-C283-intervista}



\begin{quote}
Before the new one: say the Italian for the news reporter, then the Italian for the correspondent.
\end{quote}

\subsection*{You'll want to know: l'intervista}
\textbf{l'intervista} — \textquotedblleft{}the interview\textquotedblright{}.

\begin{quote}
L'intervista è durata un'ora.
\end{quote}

\begin{quote}
\emph{The interview lasted an hour.}
\end{quote}

\emph{Fare un'intervista a qualcuno} — to interview someone.

\begin{cousinweb}
\emph{intervista} is usually said to be modelled on English \emph{interview}, which goes back to French \emph{entrevue} \textquotedblleft{}a meeting\textquotedblright{}, from \emph{entrevoir} \textquotedblleft{}to glimpse\textquotedblright{}; the root is Latin \emph{vidēre} \textquotedblleft{}to see\textquotedblright{} (Lewis \& Short).
\end{cousinweb}

\begin{grammarlens}[title={What you've built}]
One more word for telling, explaining and describing at B1.
\end{grammarlens}

\subsection*{Your turn}
\begin{itemize}
\raggedright
  \item \emph{Say it:} \emph{l'intervista}
  \item \emph{Say it:} \emph{l'intervista}, once more
  \item \emph{Say it:} a sentence of your own with \emph{l'intervista}
  \item \emph{Recall:} say the Italian for the news reporter, then the Italian for the correspondent, then say \emph{l'intervista} again
\end{itemize}

\begin{tcolorbox}[breakable,colback=teal!4,colframe=teal!35!black,title={Before you move on}]
What does l'intervista mean? (\textquotedblleft{}The interview\textquotedblright{}.) Say it once more.
\end{tcolorbox}
\section[intervistare]{intervistare — to interview}

\label{lesson:IT-C283-intervistare}



\begin{quote}
Before the new one: say the Italian for the correspondent, then the Italian for the interview.
\end{quote}

\subsection*{You'll want to know: intervistare}
\textbf{intervistare} — \textquotedblleft{}to interview\textquotedblright{}.

\begin{quote}
Ci hanno intervistato per strada.
\end{quote}

\begin{quote}
\emph{They interviewed us in the street.}
\end{quote}

Built on \emph{intervista}, the same way as in English.

\begin{cousinweb}
\emph{intervistare} is built on \emph{intervista}, which is usually traced to English \emph{interview}; the root is Latin \emph{vidēre} \textquotedblleft{}to see\textquotedblright{} (Lewis \& Short).
\end{cousinweb}

\begin{grammarlens}[title={What you've built}]
One more word for telling, explaining and describing at B1.
\end{grammarlens}

\subsection*{Your turn}
\begin{itemize}
\raggedright
  \item \emph{Say it:} \emph{intervistare}
  \item \emph{Say it:} \emph{intervistare}, once more
  \item \emph{Say it:} a sentence of your own with \emph{intervistare}
  \item \emph{Recall:} say the Italian for the correspondent, then the Italian for the interview, then say \emph{intervistare} again
\end{itemize}

\begin{tcolorbox}[breakable,colback=teal!4,colframe=teal!35!black,title={Before you move on}]
What does intervistare mean? (\textquotedblleft{}To interview\textquotedblright{}.) Say it once more.
\end{tcolorbox}
